% Generated draft: requires mathematical and source review.
\documentclass[11pt]{article}
\usepackage[margin=25mm]{geometry}
\usepackage{fontspec}
\setmainfont{Latin Modern Roman}
\usepackage{amsmath,amssymb,mathtools,unicode-math}
\setmathfont{Latin Modern Math}
\usepackage{enumitem,needspace,xurl,hyperref}
\setlength{\parindent}{0pt}
\setlength{\parskip}{5pt}
\setlength{\emergencystretch}{3em}
\newcommand{\N}{\mathbb{N}}
\newcommand{\Z}{\mathbb{Z}}
\newcommand{\Q}{\mathbb{Q}}
\newcommand{\R}{\mathbb{R}}
\newcommand{\C}{\mathbb{C}}
\newcommand{\F}{\mathbb{F}}
\newcommand{\A}{\mathbb{A}}
\newcommand{\PP}{\mathbb{P}}
\newcommand{\E}{\mathbb{E}}
\newcommand{\eps}{\varepsilon}
\newcommand{\dd}{\,\mathrm{d}}
\begin{document}
\section*{UnsolvedMath ID: 800597}
\textbf{Source problem code:} catalogue entry
\subsection{Definitions and mathematical statement}
The supplied record does not specify the ambient class of spaces or the admissible cover convention. Thus the following is a conditional formulation, with these parameters left explicit.

Let \(\mathcal K\) be the intended class of topological spaces. For each relevant space \(Y\), let \(\mathcal C(Y)\) be its family of admissible covers. For the definitions below, a cover is a family \(\mathcal U\) of proper subsets of \(Y\) such that \(\bigcup\mathcal U=Y\). If countable open covers are intended, take \(\mathcal C(Y)\) to consist of all countable such covers by open subsets. This choice, including the exclusion of \(Y\) as a cover member, requires confirmation from the source context.

Define the following subfamilies of \(\mathcal C(Y)\):
\begin{itemize}
\item \(\Omega(Y)\) consists of the covers \(\mathcal U\) such that every finite subset \(F\subseteq Y\) is contained in some \(U\in\mathcal U\).
\item \(\Gamma(Y)\) consists of the infinite covers \(\mathcal U\) such that \(\{U\in\mathcal U:y\notin U\}\) is finite for every \(y\in Y\).
\item \(\mathrm T(Y)\) consists of the covers \(\mathcal U\) such that every \(y\in Y\) belongs to infinitely many members of \(\mathcal U\), and for every \(x,y\in Y\), at least one of
\[
\{U\in\mathcal U:x\in U,\ y\notin U\},\qquad
\{U\in\mathcal U:y\in U,\ x\notin U\}
\]
is finite. These are the \(\tau\)-covers under the chosen admissibility convention.
\end{itemize}
For cover classes \(\mathcal A(Y)\) and \(\mathcal B(Y)\), the property \(\mathsf S_1(\mathcal A,\mathcal B)\) means that for every sequence \((\mathcal U_n)_{n\geq1}\) in \(\mathcal A(Y)\), there are choices \(U_n\in\mathcal U_n\) such that \(\{U_n:n\geq1\}\in\mathcal B(Y)\). The property \(\mathsf S_{\mathrm{fin}}(\mathcal A,\mathcal B)\) means that for every such sequence there are finite subfamilies \(\mathcal F_n\subseteq\mathcal U_n\), possibly empty, such that \(\bigcup_{n\geq1}\mathcal F_n\in\mathcal B(Y)\).

For a property \(P\), preservation under finite powers means
\[
\forall X\in\mathcal K\;\bigl(P(X)\Longrightarrow
\forall m\in\Z_{\geq1}\;P(X^m)\bigr),
\]
where \(X^m\) has the product topology and its cover classes are formed using the same admissibility convention. The positive integers are used here for finite powers; the supplied record does not expressly address the zeroth power.

\textbf{Problem.} Determine which, if any, of the following seven properties are preserved under finite powers:
\[
\begin{gathered}
\mathsf S_1(\Omega,\mathrm T),\qquad
\mathsf S_{\mathrm{fin}}(\Omega,\mathrm T),\\
\mathsf S_1(\mathrm T,\Gamma),\qquad
\mathsf S_1(\mathrm T,\mathrm T),\qquad
\mathsf S_1(\mathrm T,\Omega),\\
\mathsf S_{\mathrm{fin}}(\mathrm T,\mathrm T),\qquad
\mathsf S_{\mathrm{fin}}(\mathrm T,\Omega).
\end{gathered}
\]
Each property gives a separate preservation question. Specifying \(\mathcal K\) and the admissible covers from the original source context is necessary to turn this conditional formulation into an unqualified statement of the intended problem.
\subsection{Short English statement}
For each of seven selection properties involving \(\omega\)-, \(\tau\)-, and \(\gamma\)-covers, ask whether every positive finite Cartesian power of a space with that property also has it. The supplied record leaves the class of spaces and cover conventions unspecified, so the formulation retains these as unresolved parameters.
\subsection{Sources}
\begin{itemize}
\item Local UnsolvedMath registry: permanent ID \(800597\), exact source problem code \texttt{catalogue entry}, title ``Selection Principles in Mathematics --- Problem 25.'' The supplied record attributes the problem to Tsaban, with proposed year 2006, and labels its status as \texttt{open}. This is a supplied registry label, not an independently verified mathematical or literature assessment.
\item Supplied primary-source reference: Tsaban, \emph{Selection Principles and Special Sets of Reals: Open Problems} (2006), Problem 25, \url{https://arxiv.org/abs/math/0609601}. The background records extraction as \texttt{arxiv\_latex} and an access date of 2026-10-16. These are supplied provenance statements; the reference and its surrounding definitions have not been independently consulted or verified for this draft.
\item Supplied registry provenance: upstream ID \(800597\), upstream revision \texttt{d371ee917ac9647fbb06cd6fc529deec6f3dda82}; record created and updated at \texttt{2026-10-16T14:00:00Z}. These timestamps do not establish a literature-verification date or a current mathematical status.
\end{itemize}
\end{document}
